\chapter{Fundamentação teórica}
\label{cap:fundamentacao}

\section{Sentimento financeiro e empresa-alvo}

A análise de sentimentos atribui uma orientação a um texto ou a uma parte dele. No contexto desta pesquisa, são consideradas as classes positiva, negativa e neutra. O significado dessas classes depende do alvo da análise. Uma classificação do documento inteiro resume seu conteúdo dominante; uma classificação orientada à empresa procura identificar a polaridade da informação relativa à companhia indicada.

Essa distinção é relevante em notícias financeiras. Considere, como exemplo hipotético, uma matéria que relate aumento das despesas da empresa A e expansão das receitas da empresa B. Um único rótulo para a matéria não representa adequadamente as duas empresas. A unidade de análise precisa incluir o vínculo entre o texto e a entidade a que o rótulo se refere. Tarefas orientadas a empresas e ativos aparecem em \citeonline{cortis2017semeval5}, \citeonline{yuan2020target} e \citeonline{sinha2022sentfin}.

Neste trabalho, o rótulo humano expressa se o fato descrito é favorável, desfavorável ou neutro para a Sabesp, conforme as informações disponíveis no próprio texto. Esse julgamento não é uma previsão da ação. Também não utiliza o retorno posterior como evidência para escolher a classe.

É necessário distinguir neutralidade de falta de informação. Uma notícia pode ser neutra por não apresentar direcionamento identificável para a empresa; outra pode ser impossível de avaliar porque o conteúdo está incompleto. Reunir essas situações em uma única classe esconderia problemas de extração ou de pertinência do corpus. Por isso, o protocolo de continuidade prevê campos separados para a classe, a possibilidade de avaliação e a ambiguidade.

\section{Modelos de linguagem e adaptação ao domínio}

BERT é uma arquitetura de representação da linguagem que utiliza o contexto à esquerda e à direita dos termos durante o pré-treinamento e pode ser adaptada a tarefas supervisionadas, como classificação de textos (Devlin et al., 2019). Na classificação de sentimentos, uma camada de saída produz escores para as classes consideradas. O aprendizado linguístico prévio reduz a necessidade de construir manualmente todas as características textuais.

O BERTimbau disponibiliza modelos pré-treinados para o português brasileiro (Souza, Nogueira e Lotufo, 2020). Um modelo desse tipo fornece representações da linguagem, mas não deve ser confundido com um classificador de sentimentos já ajustado. Para realizar a tarefa, é necessário especificar a etapa supervisionada ou utilizar um checkpoint que já contenha a cabeça de classificação correspondente.

A adaptação ao domínio financeiro responde a diferenças de vocabulário e de uso das palavras. \citeonline{loughran2011liability} demonstram problemas na aplicação de listas de palavras de uso geral a documentos financeiros. \citeonline{araci2020finbert} investiga a adaptação de BERT à análise de sentimento financeiro em inglês. Essas contribuições justificam a atenção ao domínio, sem estabelecer que qualquer modelo financeiro seja adequado a qualquer gênero de notícia.

Na aplicação desta pesquisa, a adequação linguística, financeira e empresarial precisa ser examinada separadamente. Um modelo pode compreender português e vocabulário financeiro, mas ainda apresentar dificuldade quando a matéria contém vários alvos ou quando o julgamento esperado exige distinguir consequências favoráveis e desfavoráveis para uma companhia específica.

\section{Escores probabilísticos}

Seja uma notícia associada a uma empresa. O classificador retorna três probabilidades, representadas por $p^{+}$, $p^{-}$ e $p^{0}$, correspondentes às classes positiva, negativa e neutra. Em uma saída normalizada, a soma é igual a um. O escore utilizado no projeto é definido pela diferença

\[
d = p^{+} - p^{-}.
\]

O valor de $d$ pertence ao intervalo $[-1,1]$. Valores próximos de um indicam predominância da classe positiva; valores próximos de menos um indicam predominância da classe negativa. A classe neutra participa da distribuição de probabilidades e não é excluída do cálculo por renormalização.

Um escore próximo de zero não identifica uma única situação. Pode resultar de alta probabilidade de neutralidade ou de probabilidades semelhantes para as classes positiva e negativa. Por exemplo, as distribuições $(0{,}05;0{,}05;0{,}90)$ e $(0{,}45;0{,}45;0{,}10)$, na ordem positivo, negativo e neutro, produzem o mesmo escore zero, embora representem saídas distintas. Assim, as três probabilidades devem ser preservadas nos dados de inferência.

A maior probabilidade de saída também não deve ser interpretada automaticamente como probabilidade de acerto. Sua utilização como peso depende da calibração do classificador e da relação entre confiança e erro no corpus estudado. Uma ponderação pode amplificar previsões equivocadas se o modelo estiver excessivamente confiante.

\section{Agregação temporal e memória}

A classificação de uma notícia descreve uma observação textual. Para investigar relações ao longo do tempo, é necessário combinar observações em intervalos definidos. A contagem de matérias descreve a frequência da cobertura; a média dos escores resume sua orientação. Essas medidas têm interpretações diferentes e podem variar de forma independente.

A introdução de memória permite que observações anteriores continuem contribuindo para o estado do índice. Uma forma de representar esse comportamento é a média móvel exponencial:

\[
z_t = \alpha z_{t-1} + (1-\alpha)x_t, \qquad 0 \leq \alpha < 1,
\]

em que $x_t$ é o valor agregado no instante atual e $z_{t-1}$ é o estado anterior. Nesta convenção, valores maiores de $\alpha$ preservam uma parcela maior do passado. O parâmetro deve ser interpretado em relação à unidade temporal utilizada.

A memória introduz uma escolha de modelagem: uma observação passada perde peso de forma gradual. Isso não demonstra que o mercado processe notícias segundo a mesma função. A pertinência dessa representação deve ser avaliada por comparação com uma versão sem memória e por análise de sensibilidade do parâmetro.

Também é preciso diferenciar ausência de notícia e ausência de coleta. No primeiro caso, o procedimento pode aplicar decaimento ao estado anterior. No segundo, a cobertura do corpus é desconhecida. Tratar uma falha de coleta como dia sem notícia produziria uma alteração artificial do índice.

\section{Anotação humana e avaliação de classificação}

A anotação humana define uma referência operacional para avaliar as previsões. Sua qualidade depende da clareza das instruções e da consistência entre participantes. A concordância não torna o rótulo uma verdade econômica absoluta, mas informa se o critério pode ser aplicado de maneira relativamente estável.

A acurácia corresponde à proporção de classificações corretas. Como pode ser influenciada pela classe mais frequente, deve ser acompanhada de métricas por classe e de uma matriz de confusão. O F1 combina precisão e revocação; sua média macro atribui o mesmo peso a cada classe, tornando visíveis dificuldades que poderiam ser ocultadas pela acurácia global.

O kappa de Cohen considera a concordância observada e aquela esperada a partir das distribuições marginais de duas classificações (Cohen, 1960):

\[
\kappa = \frac{p_o-p_e}{1-p_e}.
\]

A interpretação de $\kappa$ depende da tarefa, das frequências das classes e do desenho da anotação. Não há, para esta pesquisa, um limiar universal que substitua a inspeção dos erros. \citeonline{artstein2008agreement} discutem as escolhas envolvidas na avaliação de concordância em tarefas linguísticas.

Quando há mais de dois anotadores, a medida deve corresponder ao esquema de participação. O kappa de Fleiss é uma possibilidade para avaliações nominais com desenho compatível (Fleiss, 1971). A concordância entre humanos será calculada antes da formação do consenso, para que a consolidação não apague a divergência original.

\section{Associação temporal e interpretação dos resultados}

A correlação de Pearson resume associação linear entre variáveis. A correlação de Spearman considera sua ordenação e permite examinar relações monotônicas. Nenhuma das duas demonstra, isoladamente, que uma variável cause a outra ou que forneça ganho preditivo fora da amostra.

Na pesquisa, o índice é observado antes do intervalo de retorno utilizado como desfecho. Essa separação temporal é necessária, mas insuficiente para demonstrar causalidade. Tanto as notícias quanto os preços podem responder a acontecimentos comuns. Além disso, notícias podem descrever movimentos de mercado já ocorridos.

Há ainda uma diferença entre comparar correlações e demonstrar informação incremental. Se o ITI e a média dos escores expressam componentes semelhantes da cobertura, o índice pode apresentar correlação maior sem acrescentar informação quando ambos entram em um mesmo modelo. A avaliação incremental exige verificar a contribuição do ITI condicionada às medidas de comparação, ou examinar o ganho em previsões realizadas em períodos não utilizados no ajuste.
